Un test révèle une concentration de radon 222 de \qty{\v{C}}{Bq/m^3} dans un sous-sol de \qty{\v{V}}{m^3}. La demi-vie du radon 222 est de \qty{3.82}{d}. Santé Canada recommande d'agir au-dessus de \qty{200}{Bq/m^3}.

\begin{enumerate}[label=\alph*)]
    \item Déterminez l'activité totale du radon dans le sous-sol.
    \item Combien d'atomes de radon 222 l'air du sous-sol contient-il ?
    \item On colmate la fissure par laquelle le radon entrait et on ferme le sous-sol, sans ventilation. Combien de jours faut-il pour que la concentration redescende sous \qty{200}{Bq/m^3} ?
\end{enumerate}
